Language models are increasingly relied upon for tasks requiring faithful text reproduction---copying code, quoting sources, and transcribing data. Prior work identified ``glitch tokens'' such as \texttt{SolidGoldMagikarp} that caused dramatic failures in earlier models, but it remains unclear whether these vulnerabilities persist in state-of-the-art systems, or whether new reproduction failure modes have emerged. We conduct a systematic evaluation of text reproduction fidelity across three GPT-4.1-family models using 500+ test cases spanning glitch tokens, adversarial Unicode, misspelled documents, and extreme-length texts. We find that classic glitch tokens are entirely resolved in 2025-era models---all 50 tested tokens are reproduced perfectly across all three models. Models also faithfully preserve misspelled text without auto-correcting, even at 50\% misspelling density. However, we discover a new class of ``tongue twisters'' at extreme text lengths: all models fail on long repetitive text (10k+ characters), and \gptfull paradoxically refuses to reproduce long clean text ($\geq$8k characters) while perfectly reproducing equally long misspelled text. This asymmetry suggests content-aware output filtering rather than a capability limitation. Our findings redefine the landscape of LLM reproduction failures, shifting attention from tokenizer-level artifacts to length- and content-dependent behaviors that have practical implications for any application requiring faithful text copying.
